% !TEX root = ./interim_main.tex
% 複数フォルダに分かれている場合、コンパイルするディレクトリがルートになる場合があるため必要
\documentclass[uplatex,dvipdfmx,a4j,9pt]{jsarticle}
% interim_preamble.tex

%------------------パッケージ読み込み--------------------
\usepackage{graphicx}
\usepackage{amsmath,amssymb,amsthm}
\allowdisplaybreaks
\usepackage[subrefformat=parens,labelformat=parens]{subcaption}
% \usepackage{tabularray} % platexと非互換
\usepackage{url}
\usepackage{float}
\usepackage{wrapfig}
\usepackage[top=1.5cm, bottom=1.0cm, left=1.4cm, right=1.4cm, includefoot]{geometry}

%-------------------コマンドの定義-------------------------
\newcommand{\bm}{\boldsymbol}
\newcommand{\del}{\partial}
\newcommand{\diag}{\mathrm{diag}}

\theoremstyle{definition}
\newtheorem{definition}{Definition}
\newtheorem{assumption}{Assumption}
\newtheorem{theorem}{Theorem}
\newtheorem{lemma}{Lemma}
\newtheorem{corollary}{Corollary}
\newtheorem{remark}{Remark}

\captionsetup[subfigure]{labelformat=simple}
\renewcommand{\thesubfigure}{(\alph{subfigure})}

\date{\number\year/\number\month/\number\day}

\bibliographystyle{umlab}

% ---- パッケージ群 ----------------------------------------------------------
\usepackage{multicol}
\usepackage{titlesec}   % 見出し余白調整
\usepackage{tabularray} % 表組用
\usepackage{graphicx}   % 画像挿入用
\usepackage{caption}    % キャプション調整用
\usepackage{float}      % 画像位置制御用

% ---- この部分だけ各自で編集 ----------------------------------------------
\newcommand{\JapaneseTitle}{研究題目を記入}
\newcommand{\EnglishTitle}{English Title of Your Research}
\newcommand{\ReportDate}{\today}
\newcommand{\Presenter}{氏名を記入}
\newcommand{\Supervisor}{指導教員名・職位を記入}

% ---- 見出し余白設定 --------------------------------------------------------
\titlespacing*{\section}{0pt}{1.5ex plus .2ex minus .2ex}{1ex plus .2ex}
\titlespacing*{\subsection}{0pt}{1ex plus .2ex minus .2ex}{0.8ex plus .2ex}

% ---- \maketitle 再定義 -----------------------------------------------------
\makeatletter
\def\@maketitle{%
  \let\footnote\thanks
  \begin{center}
      {\huge\bfseries \JapaneseTitle\\}
      \vspace{0.7zw}
      {\Large\bfseries \EnglishTitle}
  \end{center}
  \vspace{1.5zw}
  \hfill
  \begin{minipage}{0.45\linewidth}
      \raggedleft
      \large
      \ReportDate \\[0.8zw]
      発表者　　\Presenter　　　　 \\[0.6zw]
      指導教員　\Supervisor
  \end{minipage}
  \par\vspace{2zw}
  \setcounter{footnote}{0}%
  \let\thanks\relax\let\maketitle\relax
  \gdef\@thanks{}\gdef\@author{}\gdef\@date{}\gdef\@title{}%
}
\makeatother
\date{}

\begin{document}
\maketitle

% 本文中の「ここに〜」を各自の研究内容に置き換えてください。
% 不要な節・表・図は削除して構いません。

\section*{研究の概要と目的}
ここに研究の概要と目的を日本語で記入する。研究背景、解決したい課題、提案する方法、および期待される成果を簡潔にまとめる。

\vspace{0.5cm}
Write an English summary of the research here. Briefly describe the background, research objective, proposed approach, and expected outcomes.
\vspace{1cm}

\begin{multicols}{2}

\section{背景・先行研究調査}

\subsection{社会的背景}
ここに研究に関連する社会的背景や実用上の課題を記入する。

\subsection{学術的背景}
ここに先行研究の動向や、未解決の学術的課題を記入する。文献を引用する場合は、\cite{reference1}のように記述する。

\subsection{先行研究との位置づけ}
先行研究と比較した本研究の位置づけ、独創性、新規性、技術的意義を記入する。

\section{目的}
本研究の目的を以下に示す：
\begin{enumerate}
  \item 目的1を記入する
  \item 目的2を記入する
  \item 目的3を記入する
\end{enumerate}

\section{計画・目標}

\subsection{研究計画}
今後の研究計画、実験計画、論文執筆などの予定を記入する。

% 工程表を挿入する場合の例です。
% \begin{minipage}{\columnwidth}
%   \centering
%   \includegraphics[width=\columnwidth]{schedule-file-name}
%   \captionof{figure}{研究計画工程表}
%   \label{fig:schedule}
% \end{minipage}

\subsection{到達目標・評価項目}
本研究で目指す成果と、その評価項目を記入する。複数の目標がある場合は、箇条書きで整理する。

\section{研究手法}

\subsection{システム設計・開発}
ここに対象システム、実験装置、開発環境などを記入する。

\begin{table}[H]
\centering
\caption{システム構成（記入例）}
\begin{tblr}{colspec={X[1] X[2]},width=\linewidth}
\hline
項目 & 内容 \\
\hline
項目1 & 仕様・条件などを記入 \\
項目2 & 仕様・条件などを記入 \\
項目3 & 仕様・条件などを記入 \\
\hline
\end{tblr}
\end{table}

\subsection{提案手法}
ここに提案する手法、アルゴリズム、解析方法などを記入する。

\subsection{実験・評価方法}
評価に用いる条件、手順、比較対象、評価指標を記入する。

% 図を挿入する場合は、次の例を利用してください。
% \begin{figure}[H]
%   \centering
%   \includegraphics[width=\columnwidth]{figure-file-name}
%   \caption{図の説明}
%   \label{fig:example}
% \end{figure}

\section{まとめ}
研究の背景、目的、計画・目標、研究手法を踏まえ、現時点での進捗と今後の展望を簡潔にまとめる。

参照テスト\cite{reference1}。

\bibliography{bibfile}

\end{multicols}
\end{document}
